\section{Introduction}

A binary lift of a visible (n)-cycle carries a unique nontrivial return
receipt after one visible circuit. When the lifted successor is also the
ordered product of two involutions, a second record appears: the commutator of
the two participating operations. This paper determines exactly how those two
receipts are related.

The answer is parity-sensitive. The doubled cyclic lift itself exists at every
finite cycle length, but the commutator reaches its binary deck involution only
when the visible length is even. At odd length, the commutator and deck
involution split as independent factors. The natural encounter quotient detects
the same distinction: it erases the binary return in the even case and retains
its visible product class in the odd case.

All statements follow from three premises: two involutions (A) and (B),
their ordered product (H=AB), and the requirement that (H) have order
(2n). The relevant dihedral facts are standard \cite{dummit-foote,rotman};
the contribution here is the explicit receipt decomposition, its parity law,
and its quotient consequence.

